\chapter{Giới thiệu}
\label{chap:introduction}

\section{Bối cảnh và động lực}
\label{sec:background}

Trong hơn một thập kỷ qua, kể từ khi Satoshi Nakamoto công bố Bitcoin vào năm
2008~\cite{nakamoto2008bitcoin}, công nghệ chuỗi khối (blockchain) đã phát triển
từ một hệ thống tiền điện tử đơn lẻ thành một họ công nghệ nền tảng cho nhiều
lĩnh vực: tài chính phi tập trung, truy xuất nguồn gốc hàng hóa, quản lý tài
sản số, định danh phi tập trung và hợp đồng thông minh. Bản chất của blockchain
--- một \textit{sổ cái phân tán} (distributed ledger) bất biến, được nhiều bên
không hoàn toàn tin tưởng nhau cùng duy trì --- mang lại ba thuộc tính đặc biệt
giá trị với doanh nghiệp và nhà nước: tính minh bạch, tính
bất biến và khả năng kiểm toán mà không cần một bên trung gian
duy nhất nắm giữ quyền lực tuyệt đối.

Sự trưởng thành của công nghệ này đã được phản ánh trong định hướng chính sách
ở tầm quốc gia. Nghị quyết số 57-NQ/TW ngày 22/12/2024 của Bộ Chính trị về phát
triển khoa học, công nghệ, đổi mới sáng tạo và chuyển đổi số quốc gia đã xác
định blockchain là một trong các công nghệ chiến lược cần ưu tiên làm
chủ~\cite{nq57}. Trước đó, Quyết định số 1236/QĐ-TTg ngày 30/6/2023 đã phê duyệt
\textit{Chiến lược quốc gia về ứng dụng và phát triển công nghệ blockchain đến
năm 2030}~\cite{qd1236}, trong đó đặt ra mục tiêu xây dựng hạ tầng mạng
blockchain quốc gia ``Make in Vietnam'' với năng lực xử lý cao, độ trễ thấp và
khả năng mở rộng tới hàng trăm nút xác thực phân tán.

Một trong những yêu cầu kỹ thuật then chốt được nêu ra trong các đề án phát
triển hạ tầng blockchain quốc gia là khả năng xử lý tối thiểu 5.000 giao
dịch mỗi giây với độ trễ dưới một giây, đồng thời đảm bảo an toàn dữ liệu và
khả năng mở rộng~\cite{qd1236}. Đây không phải là những con số dễ đạt được:
chúng đòi hỏi hiểu biết sâu sắc về kiến trúc hệ phân tán, thuật toán đồng thuận,
tối ưu hiệu năng trên đường truyền mạng và tầng lưu trữ. Quan trọng hơn, để
thực sự \textit{làm chủ} công nghệ --- chứ không chỉ vận hành lại một sản phẩm
mã nguồn mở của nước ngoài --- người kỹ sư cần hiểu được từng lớp của hệ thống
hoạt động ra sao, vì sao lại thiết kế như vậy, và đâu là các điểm đánh đổi
(trade-off) cốt lõi.

Chính trong bối cảnh đó, khóa luận này đặt ra một mục tiêu học thuật rõ ràng:
\textit{tự xây dựng một nền tảng blockchain định hướng doanh nghiệp từ đầu, để
hiểu thấu đáo cơ chế bên trong của một hệ thống sổ cái phân tán có cấp phép}.
Chúng tôi không nhằm cạnh tranh với các nền tảng công nghiệp đã được kiểm chứng
như Hyperledger Fabric~\cite{androulaki2018hyperledger}, mà coi việc tái hiện
các nguyên lý thiết kế của chúng --- đặc biệt là việc tự cài đặt thuật
toán đồng thuận Raft~\cite{ongaro2014raft} --- như một con đường để làm chủ tri
thức nền tảng, một bước chuẩn bị cho việc tham gia phát triển hạ tầng blockchain
quốc gia trong tương lai.

\section{Blockchain công khai và blockchain có cấp phép}
\label{sec:public-vs-permissioned}

Không phải mọi blockchain đều giống nhau. Một sự phân loại quan trọng --- và là
nền tảng cho toàn bộ thiết kế của khóa luận --- là sự phân biệt giữa blockchain
công khai (public/permissionless) và blockchain có cấp phép
(permissioned).

Blockchain công khai như Bitcoin~\cite{nakamoto2008bitcoin} hay
Ethereum~\cite{buterin2014ethereum} cho phép bất kỳ ai cũng có thể tham gia
mạng, đọc dữ liệu, gửi giao dịch và tham gia vào quá trình đồng thuận một cách
ẩn danh. Vì các bên tham gia không có danh tính rõ ràng và có thể có hành vi
xấu, các hệ thống này phải dựa vào những cơ chế đồng thuận tốn kém như
Proof-of-Work (bằng chứng công việc) --- vốn đòi hỏi một lượng năng lượng tính
toán khổng lồ --- để chống lại tấn công Sybil và đảm bảo an toàn.

Ngược lại, blockchain có cấp phép --- nhóm mà hệ thống trong khóa luận này thuộc
về --- chỉ cho phép các thành viên \textit{được cấp phép} tham gia. Mỗi node có
một danh tính mật mã rõ ràng, thường được biểu diễn qua một cặp khóa công
khai/bí mật. Nhờ giả định về danh tính này, hệ thống có thể sử dụng các thuật
toán đồng thuận nhẹ hơn nhiều (như Raft hoặc PBFT~\cite{castro1999pbft}), đạt
throughput cao hơn và độ trễ thấp hơn đáng kể so với các blockchain công khai.
Bảng~\ref{tab:public-permissioned} tóm tắt những khác biệt cốt lõi.

\begin{table}[H]
\centering
\caption{So sánh blockchain công khai và blockchain có cấp phép}
\label{tab:public-permissioned}
\begin{tabular}{@{}p{3.2cm}p{4.2cm}p{4.2cm}@{}}
\toprule
\textbf{Tiêu chí} & \textbf{Công khai} & \textbf{Có cấp phép (dự án này)} \\
\midrule
Ví dụ điển hình & Bitcoin, Ethereum & Hyperledger Fabric, hệ thống này \\
Ai được tham gia & Bất kỳ ai & Chỉ thành viên được cấp phép \\
Định danh & Ẩn danh / bút danh & Danh tính mật mã rõ ràng \\
Cơ chế đồng thuận & Tốn kém (PoW, PoS) & Nhẹ (Raft, PBFT) \\
Throughput & Thấp (vài chục TPS) & Cao (hàng nghìn TPS) \\
Độ trễ & Cao (phút) & Thấp (giây / mili-giây) \\
Mô hình tin cậy & Không tin cậy nhau & Tin cậy một phần \\
\bottomrule
\end{tabular}
\end{table}

Mô hình có cấp phép phù hợp tự nhiên với ngữ cảnh doanh nghiệp và khu vực công:
các tổ chức tham gia (ngân hàng, cơ quan quản lý, nhà cung cấp trong chuỗi cung
ứng) đều biết danh tính của nhau, chịu sự ràng buộc pháp lý, và cần một hệ thống
vừa minh bạch vừa có hiệu năng đủ cao để phục vụ nghiệp vụ thực tế. Đây chính là
phân khúc mà Hyperledger Fabric nhắm tới, và cũng là định hướng của khóa luận.

\section{Phát biểu bài toán}
\label{sec:problem}

Xây dựng một nền tảng blockchain doanh nghiệp hoàn chỉnh từ đầu đặt ra một loạt
thách thức kỹ thuật đan xen nhau, mà nếu chỉ sử dụng lại một nền tảng có sẵn thì
người học sẽ không bao giờ phải đối mặt một cách trực tiếp:

\begin{enumerate}
  \item \textbf{Bài toán đồng thuận.} Khi nhiều node cùng nhận giao dịch gần
  như đồng thời, làm thế nào để tất cả thống nhất về một \textit{thứ tự toàn
  phần} (total order) duy nhất của các giao dịch? Đây là bài toán đồng thuận
  kinh điển trong hệ phân tán, với kết quả bất khả thi nổi tiếng
  FLP~\cite{fischer1985flp} cho thấy không thể có thuật toán đồng thuận vừa an
  toàn vừa luôn kết thúc trong mạng bất đồng bộ có lỗi. Việc tự cài đặt một
  thuật toán đồng thuận đúng đắn (ở đây là Raft) đòi hỏi xử lý cẩn thận các tình
  huống bầu cử, mất kết nối, và phân ly mạng (split-brain).

  \item \textbf{Bài toán thực thi hợp đồng thông minh an toàn.} Hợp đồng thông
  minh phải chạy một cách \textit{xác định} (deterministic) --- cùng đầu vào
  luôn cho cùng kết quả trên mọi node --- và \textit{cô lập} (sandboxed) để
  không thể phá hoại hệ thống chủ. Làm thế nào để hỗ trợ nhiều ngôn ngữ lập
  trình mà vẫn đảm bảo hai thuộc tính này?

  \item \textbf{Bài toán xác thực và toàn vẹn.} Làm thế nào để xác thực ``ai
  gửi giao dịch'' mà không cần mật khẩu, và đảm bảo dữ liệu sổ cái không thể bị
  sửa đổi một cách âm thầm?

  \item \textbf{Bài toán hiệu năng.} Làm thế nào để đạt throughput hàng nghìn
  giao dịch mỗi giây trong khi vẫn giữ độ trễ ở mức chấp nhận được? Đâu là các
  nút nghẽn cổ chai (bottleneck) trên đường truyền, và làm thế nào để nhận diện
  cũng như loại bỏ chúng?

  \item \textbf{Bài toán giao tiếp mạng.} Làm thế nào để các thành phần phân tán
  giao tiếp một cách an toàn, có định danh mật mã, và chịu được việc node tham
  gia/rời mạng động?
\end{enumerate}

Mỗi thách thức trên tương ứng với một quyết định thiết kế cụ thể trong hệ thống.
Khóa luận này đặt mục tiêu \textit{trình bày và lý giải} từng quyết định đó, dựa
trên mã nguồn thực tế đã được hiện thực và kiểm thử.

\section{Mục tiêu và phạm vi}
\label{sec:objectives-scope}

\subsection{Mục tiêu}

Mục tiêu tổng quát của khóa luận là thiết kế, hiện thực và đánh giá một
nền tảng blockchain có cấp phép định hướng doanh nghiệp, tự xây dựng từ đầu.
Mục tiêu này được cụ thể hóa thành các mục tiêu thành phần:

\begin{itemize}
  \item Tự cài đặt thuật toán đồng thuận Raft từ đầu (không dùng thư
  viện), bao gồm bầu lãnh đạo, nhân bản log, cắt block và đồng bộ --- với một
  biến thể bầu lãnh đạo theo độ ưu tiên.

  \item Xây dựng một máy ảo hợp đồng thông minh dựa trên WebAssembly,
  cho phép viết hợp đồng bằng ngôn ngữ bậc cao và thực thi an toàn, xác định.

  \item Hiện thực mô hình xác nhận giao dịch Execute--Order--Validate
  theo triết lý của Hyperledger Fabric, với chữ ký số Ed25519 và kiểm tra
  endorsement.

  \item Xây dựng tầng lưu trữ bất biến (append-only) cho chuỗi block và world
  state UTXO trên LevelDB với cập nhật nguyên tử.

  \item Xây dựng giao diện web (Explorer) hiển thị dữ liệu theo thời gian thực.

  \item Thiết lập phương pháp đo lường hiệu năng và đánh giá hệ thống dưới tải
  cao, hướng tới mục tiêu 5.000 TPS.
\end{itemize}

\subsection{Phạm vi}

Để đảm bảo chiều sâu, khóa luận tập trung vào luồng xử lý cốt lõi của
một giao dịch từ lúc được gửi đến lúc được ghi vĩnh viễn. Cụ thể, báo cáo đi sâu
vào bốn thành phần: Core Service, Ordering Service (trọng tâm), Committing Peer
và Blockchain Explorer, cùng cơ sở dữ liệu PostgreSQL đóng vai trò mirror.

Một số thành phần phụ trợ phục vụ vận hành --- như công cụ điều phối cụm
(orchestrator) và giao diện giám sát cụm trong thư mục \texttt{orderingservice/}
--- nằm ngoài phạm vi trình bày, vì chúng là công cụ quản trị, không
thuộc luồng xử lý giao dịch chính. Tương tự, khóa luận tập trung vào mô hình tin
cậy của blockchain có cấp phép (chống lỗi dừng --- crash fault tolerance), và
chỉ \textit{thảo luận} (chứ không hiện thực đầy đủ) khả năng chịu lỗi Byzantine.

\section{Đóng góp của khóa luận}
\label{sec:contributions}

Những đóng góp chính của khóa luận bao gồm:

\begin{enumerate}
  \item \textbf{Một hiện thực Raft hoàn chỉnh, tự viết}, với biến thể bầu lãnh
  đạo theo độ ưu tiên (priority-based leader election) thay cho cơ chế timeout
  ngẫu nhiên của Raft chuẩn. Biến thể này loại bỏ hiện tượng chia phiếu (split
  vote) và cho phép hội tụ xác định, đồng thời vẫn giữ nguyên tắc an toàn dựa
  trên quorum đa số trên tổng số thành viên để chống split-brain.

  \item \textbf{Một kiến trúc tách tầng Execute--Order--Validate} được hiện thực
  bằng bốn tiến trình độc lập giao tiếp qua \texttt{libp2p}, minh họa cụ thể
  cách triết lý thiết kế của Hyperledger Fabric vận hành trong thực tế.

  \item \textbf{Một máy ảo hợp đồng thông minh WebAssembly} với cơ chế cache
  module đã biên dịch và pool sandbox dựng sẵn, đạt khả năng thực thi song song
  hiệu năng cao.

  \item \textbf{Một loạt tối ưu hiệu năng trên đường nóng} (được đánh số
  OPT-1 đến OPT-8), giúp tăng số giao dịch trên mỗi block từ khoảng 71 lên hơn
  1000 dưới tải bền vững, đưa hệ thống tiệm cận mục tiêu 5.000 TPS.

  \item \textbf{Một phương pháp đo lường hiệu năng ba lớp} (k6, submit, commit)
  cho phép phân biệt rõ ràng độ trễ chấp nhận của cổng vào với độ trễ
  end-to-end, và giải thích hiện tượng backlog dưới tải bão hòa thông qua định
  luật Little~\cite{little1961law}.
\end{enumerate}

\section{Phương pháp tiếp cận}
\label{sec:approach}

Khóa luận áp dụng phương pháp \textit{xây dựng để hiểu} (build-to-understand):
thay vì chỉ phân tích lý thuyết, chúng tôi hiện thực toàn bộ hệ thống bằng ngôn
ngữ Go cho ba dịch vụ backend và React cho giao diện, sau đó tiến hành kiểm thử
chức năng và đo lường hiệu năng thực nghiệm. Mỗi quyết định thiết kế đều được
đối chiếu với (a)~nguyên lý lý thuyết tương ứng (bài báo Raft, kiến trúc Fabric,
đặc tả WebAssembly...) và (b)~kết quả đo lường thực tế. Cách tiếp cận này cho
phép phát hiện và lý giải các điểm đánh đổi mà chỉ khi tự hiện thực mới bộc lộ
ra, ví dụ như nghẽn cổ chai do tranh chấp khóa khi ghi log theo từng giao dịch,
hay sự khác biệt giữa độ trễ HTTP và độ trễ ghi sổ thực sự.

\section{Kết quả đạt được}
\label{sec:results-summary}

Hệ thống đã được hiện thực hoàn chỉnh và chạy được luồng giao dịch
end-to-end. Trong các phép đo tải tham khảo (gửi 5.000 yêu cầu mỗi giây trong
60 giây), hệ thống đạt:

\begin{itemize}
  \item Tốc độ chấp nhận (submit) bền vững khoảng 4.691 giao dịch/giây, đỉnh
  khoảng 5.017 giao dịch/giây.
  \item Tốc độ ghi sổ (commit) bền vững khoảng 3.921 giao dịch/giây, đỉnh khoảng
  6.000 giao dịch/giây.
  \item Trung bình khoảng 989 giao dịch trên mỗi block (gần đúng kích thước lô
  mục tiêu 1.000).
\end{itemize}

Các kết quả này, cùng với phân tích hiện tượng backlog khi tốc độ gửi vượt tốc
độ ghi, được trình bày chi tiết ở chương đánh giá. Chúng cũng là cơ sở để đề
xuất các hướng cải thiện trong các chương sau.

\section{Cấu trúc khóa luận}
\label{sec:thesis-structure}

Phần còn lại của khóa luận được tổ chức như sau:

\begin{itemize}
  \item \textbf{Chương~\ref{chap:related-work} --- Các công trình liên quan:}
  trình bày nền tảng lý thuyết và các công trình tiền đề: Bitcoin và mô hình
  UTXO, Ethereum và hợp đồng thông minh, Hyperledger Fabric và mô hình
  Execute--Order--Validate, các thuật toán đồng thuận (Paxos, Raft, PBFT),
  WebAssembly cho hợp đồng thông minh, và định vị của khóa luận so với các công
  trình này.

  \item \textbf{Chương~\ref{chap:methodology} --- Phương pháp và Thiết kế hệ
  thống:} trình bày chi tiết kiến trúc tổng thể, mô hình dữ liệu, và thiết kế
  của từng thành phần (Core Service, Ordering Service, Committing Peer,
  Explorer), cơ sở dữ liệu, cùng phương pháp đo lường hiệu năng.

  \item Các chương tiếp theo (sẽ được hoàn thiện) trình bày hiện thực chi tiết,
  kết quả thực nghiệm, đánh giá và các đề xuất cải thiện về tốc độ, lưu trữ và
  bảo mật.
\end{itemize}
